% Options for packages loaded elsewhere
\PassOptionsToPackage{unicode}{hyperref}
\PassOptionsToPackage{hyphens}{url}
%
\documentclass[
]{article}
\usepackage{amsmath,amssymb}
\usepackage{iftex}
\ifPDFTeX
  \usepackage[T1]{fontenc}
  \usepackage[utf8]{inputenc}
  \usepackage{textcomp} % provide euro and other symbols
\else % if luatex or xetex
  \usepackage{unicode-math} % this also loads fontspec
  \defaultfontfeatures{Scale=MatchLowercase}
  \defaultfontfeatures[\rmfamily]{Ligatures=TeX,Scale=1}
\fi
\usepackage{lmodern}
\ifPDFTeX\else
  % xetex/luatex font selection
  \setmainfont[]{DejaVu Serif}
  \setmonofont[]{DejaVu Sans Mono}
\fi
% Use upquote if available, for straight quotes in verbatim environments
\IfFileExists{upquote.sty}{\usepackage{upquote}}{}
\IfFileExists{microtype.sty}{% use microtype if available
  \usepackage[]{microtype}
  \UseMicrotypeSet[protrusion]{basicmath} % disable protrusion for tt fonts
}{}
\makeatletter
\@ifundefined{KOMAClassName}{% if non-KOMA class
  \IfFileExists{parskip.sty}{%
    \usepackage{parskip}
  }{% else
    \setlength{\parindent}{0pt}
    \setlength{\parskip}{6pt plus 2pt minus 1pt}}
}{% if KOMA class
  \KOMAoptions{parskip=half}}
\makeatother
\usepackage{xcolor}
\usepackage[margin=0.72in]{geometry}
\usepackage{longtable,booktabs,array}
\usepackage{calc} % for calculating minipage widths
% Correct order of tables after \paragraph or \subparagraph
\usepackage{etoolbox}
\makeatletter
\patchcmd\longtable{\par}{\if@noskipsec\mbox{}\fi\par}{}{}
\makeatother
% Allow footnotes in longtable head/foot
\IfFileExists{footnotehyper.sty}{\usepackage{footnotehyper}}{\usepackage{footnote}}
\makesavenoteenv{longtable}
\setlength{\emergencystretch}{3em} % prevent overfull lines
\providecommand{\tightlist}{%
  \setlength{\itemsep}{0pt}\setlength{\parskip}{0pt}}
\setcounter{secnumdepth}{-\maxdimen} % remove section numbering
\ifLuaTeX
  \usepackage{selnolig}  % disable illegal ligatures
\fi
\usepackage{bookmark}
\IfFileExists{xurl.sty}{\usepackage{xurl}}{} % add URL line breaks if available
\urlstyle{same}
\hypersetup{
  hidelinks,
  pdfcreator={LaTeX via pandoc}}

\author{}
\date{}

\begin{document}

{
\setcounter{tocdepth}{3}
\tableofcontents
}
\section{EX12\_NEXT\_0 --- Rapport scientifique final et audit de
complétude}\label{ex12_next_0-rapport-scientifique-final-et-audit-de-compluxe9tude}

\textbf{Date :} 4 novembre 2026\\
\textbf{Version du protocole exécuté :} \texttt{prompt\_ex12.json}
v1.0.1 (amendement de capacité P\_trainable)\\
\textbf{Rapport scientifique directeur :}
\texttt{rapport\_prigramme\_expermentale\_EX12\_NEXT\_0.tex}\\
\textbf{Workspace :} \texttt{EX12\_NEXT\_0/}\\
\textbf{Périmètre :} phases P0 à P7, sans extension post hoc des
holdouts.

\subsection{Résumé exécutif}\label{ruxe9sumuxe9-exuxe9cutif}

EX12 a construit et évalué une famille HC-RBF (Hybrid Coverage-Adaptive
Radial Basis Function) en comparaison avec des MLP appariés en nombre de
paramètres entraînables. Le critère historique GE4X a été abandonné : le
comparateur principal est à P\_trainable ±1 \%, et les paramètres
stockés non entraînables, buffers, normaliseurs et octets de checkpoint
sont reportés séparément.

Les phases P0--P7 ont chacune une porte PASS. Cela atteste la complétude
et l'intégrité des campagnes selon leurs critères respectifs ; cela ne
signifie pas que HC-RBF a gagné scientifiquement. Le verdict
expérimental est \textbf{mixte sur la qualité prédictive, défavorable
sur l'efficacité de bout en bout et faible en extrapolation}.

Sur les quatre tâches confirmatoires synthétiques P7, 11/16 contrastes
primaires passent la non-infériorité ajustée (marge NMAE +0,02), et 7/16
rejettent l'égalité à zéro après correction bilatérale de Holm. Les
résultats favorables dépendent de la tâche ; sur \texttt{f1\_sin}, les
quatre variantes HC-RBF sont battues en interpolation par le MLP
apparié. En extrapolation, l'écart moyen candidat−MLP est +0,20559 NMAE
sur les 16 contrastes et seulement 4/16 favorisent HC-RBF. Les
benchmarks P5 canoniques montrent une latence batch 1 p50 de 1,82× à
2,25× le MLP sur le pipeline complet, un temps d'entraînement de
pipeline enregistré de 1,44× à 1,62×, et une taille exacte des tenseurs
de 1,00× pour R0/C2 et 1,27--1,28× pour C3/C4. Aucune variante ne
satisfait donc les objectifs prioritaires de vitesse et de ressources
sur le matériel mesuré.

La présente conclusion n'établit ni l'impossibilité d'améliorer HC-RBF,
ni une supériorité générale du MLP pour tous problèmes. Elle décrit
exactement le périmètre testé dans EX12.

\subsection{1. Gouvernance, sources et
protocole}\label{gouvernance-sources-et-protocole}

\subsubsection{1.1 Sources historiques
consultées}\label{sources-historiques-consultuxe9es}

L'archive \texttt{last\_sources\_3.zip} a été vérifiée et extraite
récursivement dans \texttt{/mnt/data/last\_sources\_3} après
rétablissement de l'accès par \texttt{/bin/sh}. Le processus avait
\texttt{uid=0(root)}. L'extraction principale et les archives imbriquées
ont passé les contrôles d'intégrité ZIP consignés dans P0.

Les prompts EX3, EX4, EX5, EX6, EX7, EX8 et EX11 ont été lus et audités.
\textbf{Aucun \texttt{prompt\_ex10.json} n'était présent} : les règles
EX10 ont été reconstruuites à partir de son README, de ses rapports,
scripts et audits, sans présenter ces pièces comme un prompt. EX7
réutilise le corps EX6 avec des addenda ; EX8 a deux jeux de métadonnées
(v1.0 et v1.1) avec corps identique ; EX11 est identifié comme une
reconstruction versionnée. \texttt{arch\_up\_test\_0.zip} est absent de
l'archive ; \texttt{arch\_up\_test\_1.zip} est présent et a été audité.
Ces absences ne sont pas comblées par des résultats inventés.

\subsubsection{1.2 Règles rouges et
équité}\label{ruxe8gles-rouges-et-uxe9quituxe9}

Le protocole interdit : d'inférer une réussite ou un échec avant les
résultats ; de manipuler données, splits, seeds, métriques, modèles ou
temps ; de supprimer les résultats défavorables ; de sélectionner sur
test/holdout ; d'ajuster les prétraitements sur validation/test ; de
relancer sélectivement des seeds ; de présenter un microbenchmark
théorique comme une mesure matérielle ; de confondre inspectabilité et
interprétabilité ; et de mélanger les anciens artefacts EX3--EX11 aux
données EX12.

Les runs et registres sont append-only. Les échecs initiaux,
interruptions, reprises, essais d'auditeur échoués et lacunes
documentaires sont conservés et explicitement signalés.

\subsubsection{1.3 Changement de critère de
capacité}\label{changement-de-crituxe8re-de-capacituxe9}

P0 a initialement enregistré une règle \texttt{P\_total} héritée.
L'analyse de conformité P4 a identifié que le rapport directeur exigeait
en réalité un appariement de \textbf{paramètres entraînables} à ±1 \%.
La correction a été pré-enregistrée avant P4 additionnel et avant tout
test P7 dans \texttt{PROMPT\_CHANGE\_CONTROL\_v1.0.1.json}, et intégrée
dans \texttt{prompt\_ex12.json} v1.0.1. Le gel P7 vérifie 320 paires et
un écart relatif maximal de 0,9238 \%. Le GE4X historique n'est pas un
objectif EX12.

\subsection{2. Architecture et variantes réellement
implémentées}\label{architecture-et-variantes-ruxe9ellement-impluxe9mentuxe9es}

Pour une entrée normalisée, l'encodeur partagé calcule
\texttt{z(x)=tanh(W\_e\ x+b\_e)}. Les centres sont divisés en centres
globaux fixes et centres locaux adaptatifs. Les largeurs diagonales
positives sont bornées ; chaque centre produit une activation gaussienne
\texttt{phi\_k(x)=exp(-d\_k²/2)} ; un readout radial combine les
activations. C4 ajoute un chemin affine partagé sur la représentation
latente. Aucun MLP caché n'est dissimulé dans le readout HC-RBF.

\begin{longtable}[]{@{}
  >{\raggedright\arraybackslash}p{(\columnwidth - 4\tabcolsep) * \real{0.3333}}
  >{\raggedright\arraybackslash}p{(\columnwidth - 4\tabcolsep) * \real{0.3333}}
  >{\raggedright\arraybackslash}p{(\columnwidth - 4\tabcolsep) * \real{0.3333}}@{}}
\toprule\noalign{}
\begin{minipage}[b]{\linewidth}\raggedright
Code
\end{minipage} & \begin{minipage}[b]{\linewidth}\raggedright
Implémentation / rôle
\end{minipage} & \begin{minipage}[b]{\linewidth}\raggedright
Statut expérimental
\end{minipage} \\
\midrule\noalign{}
\endhead
\bottomrule\noalign{}
\endlastfoot
\texttt{M0\_CORE} & MLP de référence apparié à R0/C2 dans les premières
phases & Implémenté, entraîné P2/P3/P5 \\
\texttt{M0\_AFFINE} & MLP de référence historique apparié à C4 &
Implémenté, entraîné P2/P3/P5 \\
\texttt{M0\_GRAPH} & MLP de contrôle du bras graphique & Implémenté et
exploré en P2 \\
\texttt{R0} & RBF adaptative, centres locaux KMeans train-only &
Implémenté, exploré puis candidat P3--P7 \\
\texttt{C1} & Centres globaux fixes uniquement & Implémenté et exploré
en P2 ; non promu en P3 \\
\texttt{C2} & Centres locaux initialisés par sélection résiduelle
diversifiée train-only & Implémenté, exploré puis candidat P3--P7 \\
\texttt{C3} & Couverture hybride centres globaux fixes + centres locaux
adaptatifs & Implémenté, exploré puis candidat P3--P7 \\
\texttt{C4} & C3 + correction affine partagée & Implémenté, exploré puis
candidat P3--P7 \\
\texttt{G1} & Graphe réel de coordination conditionnée par l'entrée &
Implémenté et exploré en P2, non promu dans la campagne confirmatoire \\
\texttt{G2} & Contrôle graphe randomisé à séquence de degrés préservée &
Implémenté et exploré en P2, non promu dans la campagne confirmatoire \\
\texttt{C5} & Modulation faible rang & Différé ; pas de variante P2
complète exécutée \\
\texttt{C6} & Sélection top-q parcimonieuse & Différé ; pas de variante
P2 complète exécutée \\
\end{longtable}

Le banc P1 a étudié des prototypes de kernels de distance, top-q, faible
rang et messages dynamiques. Les kernels top-q, faible rang et messages
dynamiques n'ont pas été promus dans le chemin de base : leur coût
mesuré n'apportait pas de gain end-to-end fiable. Les prototypes P1 ne
doivent pas être confondus avec des variantes C5/C6 entraînées et
évaluées en P2/P3.

\subsection{3. Résumé des phases et portes de
validation}\label{ruxe9sumuxe9-des-phases-et-portes-de-validation}

\begin{itemize}
\tightlist
\item
  \textbf{P0 --- PASS, 17/17 :} inventaire, audit des prompts, intégrité
  ZIP, pré-enregistrement, graines et règles rouges.
\item
  \textbf{P1 --- PASS, 21/21 :} équivalence numérique des
  sorties/gradients et microbenchmarks CPU des kernels.
\item
  \textbf{P2 --- PASS, 17/17 :} criblage exploratoire train/validation,
  80 configurations logiques et 10 variantes P2.
\item
  \textbf{P3 --- PASS, 19/19 :} 480 runs exploratoires multi-graines sur
  huit tâches, sans test/holdout.
\item
  \textbf{P4 --- PASS, 11/11 :} 320 paires appariées en paramètres
  entraînables à ±1 \%, avec 160 MLP contrôles additionnels C3/C4.
\item
  \textbf{P5 --- PASS, 15/15 :} benchmark ressources V3 : 560 modèles,
  13 440 cellules, 1 075 200 latences.
\item
  \textbf{P6 --- PASS, 24/24 :} ablations, permutations, diversité,
  stabilité inter-graines et falsification sur 320 modèles.
\item
  \textbf{P7 --- PASS, 12/12 :} gel de 640 modèles, génération des tests
  après gel, scoring unique et audit confirmatoire.
\end{itemize}

Le PASS de chaque phase signifie que les critères de complétude et
d'intégrité de cette phase ont été vérifiés ; le PASS n'est jamais
interprété comme un résultat prédictif positif.

\subsection{4. Résultats détaillés par
phase}\label{ruxe9sultats-duxe9tailluxe9s-par-phase}

\subsubsection{P1 --- kernels et équivalence
informatique}\label{p1-kernels-et-uxe9quivalence-informatique}

Deux microbenchmarks canoniques indépendants (RUN003/RUN004) ont chacun
passé 126 cas numériques et 24 000 mesures brutes. L'erreur maximale de
sortie complète est 0 ; l'erreur maximale de gradient est
\texttt{2,98×10⁻⁸}, sous tolérance \texttt{1e-6}. Le broadcasting a une
accélération médiane de 8,46× et 8,89× contre la référence à boucles ;
l'accélération du chemin HC complet est 4,16× et 4,25× sur CPU.

Le top-q a été plus lent que le calcul dense dans toutes les
comparaisons appariées du protocole ; le faible rang et les messages à
trois tours présentent également un coût élevé. La vectorisation
améliore le noyau mais ne garantit pas une accélération de
l'architecture complète.

\textbf{Lacune de reproductibilité :} les fichiers de résultats
P1\_RUN\_001 et P1\_RUN\_002 sont conservés, mais leurs scripts source
exacts ne sont pas récupérés. RUN001 est un prototype numérique échoué ;
RUN002 est exclu des comparaisons top-q non équivalentes. RUN003/RUN004
ont scripts et hashes exacts. Voir
\texttt{RAPPORT\_LACUNES\_REPRODUCTIBILITE\_P1\_RUN001\_RUN002.md}.

\subsubsection{P2 --- criblage exploratoire
contrôlé}\label{p2-criblage-exploratoire-contruxf4luxe9}

Le registre contient 80 succès logiques, 80 premières tentatives
échouées (erreur \texttt{yva} manquante), et six interruptions de
reprise ; ces échecs et interruptions restent conservés. Dix variantes
P2 ont été réellement implémentées/explorées : trois baselines MLP
(CORE/AFFINE/GRAPH), R0, C1, C2, C3, C4, G1 et G2. C5/C6 restent
\texttt{DEFERRED\_NOT\_IMPLEMENTED}.

Le criblage a retenu R0/C2/C3/C4 pour P3. Ce résultat de deux seeds et
quatre fonctions est exploratoire, pas confirmatoire.

\subsubsection{P3 --- tâches et généralisation
exploratoires}\label{p3-tuxe2ches-et-guxe9nuxe9ralisation-exploratoires}

480 runs logiques sur huit tâches et dix seeds : quatre fonctions
principales (\texttt{f1} à \texttt{f4}), deux compositions diagnostiques
(\texttt{g5},\texttt{g6}), Diabetes tabulaire et classification binaire
synthétique. Six slots de modèles : R0/C2/C3/C4 et deux MLP initiaux. 71
fichiers/664 arrays de données contrôlés ; 1 620 métriques recalculées ;
13 premières tentatives d'audit échouées conservées, sans interruption
non résolue.

Les NMAE moyennes d'interpolation sur dix seeds montrent des résultats
dépendant de la fonction : \texttt{f1\_sin} défavorise toutes les sondes
(MLP 0,1050 contre HC-RBF 0,1741--0,2792) ; \texttt{f2\_mul} favorise
HC-RBF (MLP 0,0700, HC-RBF 0,0203--0,0327) ; \texttt{f3\_square\_plus}
favorise certaines sondes (MLP 0,0767, C4 0,0367) ;
\texttt{f4\_sin\_product} favorise R0/C2 et certains cas. Les tâches
\texttt{g5/g6} demeurent diagnostiques.

Sur classification binaire synthétique, les log-loss HC-RBF restent
supérieures aux MLP appariés : deltas moyens +0,0197 (R0), +0,0116 (C2),
+0,0187 (C3), +0,0328 (C4). Sur Diabetes validation, les NMAE HC-RBF
sont légèrement inférieures dans ce split unique ; le holdout de 67
observations reste scellé.

\subsubsection{P4 --- appariement MLP}\label{p4-appariement-mlp}

L'audit a détecté que le \texttt{P\_total} pouvait masquer une
différence de paramètres entraînables. 160 MLP contrôles supplémentaires
ont été entraînés pour C3/C4 ; 320 paires candidates--MLP passent la
contrainte \texttt{P\_trainable\ ±1\ \%}. Les 8 960 valeurs de métriques
et 2 800 empreintes de fichiers sont recalculées/vérifiées sans
divergence. Le test/holdout est resté inaccessible. Les anciennes
synthèses \texttt{P\_total} sont conservées comme analyses historiques
mais supersédées pour l'équité principale.

\subsubsection{P5 --- ressources
canoniques}\label{p5-ressources-canoniques}

Le gate P5 canonique est basé sur 560 modèles, 13 440 cellules de
timing, 1 075 200 latences élémentaires, 2 560 lignes de latence
appariées, 320 paires mémoire et 320 paires de temps de pipeline. CPU
mono-thread, float32, batches 1/16/64/256, deux chemins, 3 blocs
indépendants de 80 mesures. Les ratios p50 pipeline batch 1 candidat/MLP
sont R0 1,825×, C2 1,824×, C3 1,955×, C4 2,254× ; au batch 256 ils sont
2,586×, 2,609×, 2,922×, 3,104×. Fraction de paires plus rapides au batch
1 : 0 pour les quatre variantes.

Les ratios de temps de pipeline enregistrés sont R0 1,483×, C2 1,515×,
C3 1,445×, C4 1,622×. Ce temps inclut initialisation, validation,
diagnostics et écritures ; ce n'est pas un chronométrage pur de
l'optimiseur. Les rapports de stockage exact des tenseurs sont R0
0,996×, C2 0,996×, C3 1,269×, C4 1,280×. Aucun GPU ni compteur d'énergie
n'était disponible ; RSS est process-level, non attribuable à un modèle
isolé.

\textbf{Lacune de reproductibilité :} le script V2 exact référencé par
13 416 cellules historiques n'a pas été retrouvé. Les données sont
conservées ; une reproduction V3 indépendante et hashée couvre 13 440
cellules et est comparée aux anciennes. Voir
\texttt{RAPPORT\_LACUNE\_REPRODUCTIBILITE\_P5\_V2.md} et
\texttt{P5\_RESOURCE\_REPORT\_V3.md}.

\textbf{Erratum de comptage :} une synthèse d'avancement avait attribué
640 modèles/256 000 latences à P5. C'était une confusion avec des
artefacts P7 ; les comptes canoniques P5 sont 560 modèles et 1 075 200
latences. Voir \texttt{10\_reports/P5/P5\_COUNT\_RECONCILIATION.md}.

\subsubsection{P6 --- complémentarité et
spécialisation}\label{p6-compluxe9mentarituxe9-et-spuxe9cialisation}

320 modèles candidats, 10 240 unités modèle-centre, 34 560 lignes
d'ablation, 1 600 permutations invariantes, 1 440 paires modèles
cross-seed, 46 080 appariements centre-centre et 1 024 hypothèses de
rôle. Aucune divergence d'intégrité P3/P6 détectée. Les permutations
sont group-preserving et respectent les tolérances numériques
pré-enregistrées.

Résultat scientifique : 549/1 024 centres passent le seuil de stabilité,
9/1 024 passent le test d'utilité après Holm, et seulement \textbf{7/1
024} passent les deux. Les sept sont limités à
\texttt{g5\_sin\_mul\_plus\_square} : C2=3, C3=2, C4=2, R0=0. La
spécialisation utile et stable n'est donc pas une propriété générale de
l'architecture.

\textbf{Lacune de traçabilité :} 123 terminal records P6 v1 ne
contiennent pas le hash du script exécuté, même si la source v1 existe
actuellement. Le lien cryptographique exact au moment de ces 123
exécutions ne peut pas être prouvé depuis les lignes terminales. Une NPZ
partielle d'une tentative interrompue est conservée comme artefact
orphelin attendu.

\textbf{Incident \texttt{manifest\_map} :} le rapport séparé
\texttt{P6\_MANIFEST\_MAP\_INCIDENT\_REPORT.pdf} et son dossier d'audit
montrent que le manifeste persistant \texttt{P3\_DATA\_MANIFEST.json}
est présent (71 entrées), que les runners P6 contiennent la fonction
\texttt{load\_manifest\_map()}, et que 320/320 sources se rapprochent
sans fichier manquant ni divergence SHA-256. L'audit n'a trouvé aucune
suppression persistante ou altération des données. L'incident ne
constitue pas une preuve de sabotage.

\subsubsection{P7 --- confirmation finale}\label{p7-confirmation-finale}

P7 a gelé 640 modèles (4 tâches primaires × 20 graines × 8 slots), 320
paires candidates--MLP, 4 257 artefacts pré-test et les scripts. La
génération test a eu lieu après gel : 80 fichiers, 640 arrays, quatre
régimes et 1 000 observations par régime. Scoring unique : 640 modèles ×
4 régimes = 2 560 évaluations ; audit 11/11, aucune divergence et aucune
sélection/réentraînement après le test.

La métrique primaire est NMAE normalisée par l'écart-type de cible
calculé sur train. Sur les 16 contrastes d'interpolation, 11 passent la
non-infériorité ajustée (marge +0,02), et 7 montrent une différence
significative après Holm bilatéral. Trois contrastes significatifs
défavorisent les sondes sur \texttt{f1\_sin} ; quatre favorisent des
variantes sur \texttt{f3\_square\_plus} et \texttt{f4\_sin\_product}.
Aucune variante ne gagne universellement.

Par régime, les différences moyennes NMAE candidat−MLP sont
interpolation −0,00571, frontière −0,03192, extrapolation +0,20559,
shift +0,00298. L'extrapolation est fortement défavorable sur
\texttt{f1\_sin}, \texttt{f2\_mul}, \texttt{f3\_square\_plus} (toutes
les variantes sont plus mauvaises que leur MLP) ; seul
\texttt{f4\_sin\_product} est favorable aux sondes en moyenne. La
frontière favorise HC-RBF dans 15/16 contrastes. Le shift est variable
par tâche.

Le rapport détaillé P7 et tous les tableaux par seed sont dans
\texttt{10\_reports/P7/P7\_CONFIRMATORY\_FINAL\_REPORT.md} et
\texttt{09\_analysis/P7/}. Aucun test confirmatoire de Diabetes ou de
classification n'a été réalisé dans P7 ; leurs résultats P3/P5 ne sont
pas promus en résultats de test. Le protocole P7 primaire porte sur les
quatre fonctions synthétiques déclarées.

\subsection{5. Lacunes documentaires et
reproductibilité}\label{lacunes-documentaires-et-reproductibilituxe9}

Les lacunes suivantes sont explicitement conservées, sans conclusion non
démontrée de sabotage ou de vol :

\begin{enumerate}
\def\labelenumi{\arabic{enumi}.}
\tightlist
\item
  \texttt{prompt\_ex10.json} absent ; protocole EX10 reconstruit depuis
  les artefacts EX10 disponibles.
\item
  \texttt{arch\_up\_test\_0.zip} absent ; aucune donnée ni aucun
  résultat EX0 inféré.
\item
  Scripts sources exacts P1\_RUN\_001/002 absents ; résultats et logs
  conservés, mais reproduction byte-identique non démontrable.
  RUN003/004 sont canoniques et reproductibles.
\item
  Script P5 V2 exact absent pour 13 416 cellules ; reproduction
  indépendante V3 complète, hashée et comparée. Les 13 440 cellules V3
  sont reproductibles.
\item
  123 lignes terminales P6 v1 sans hash de script ; source v1 actuelle
  disponible, mais identité historique des octets non attestée
  cryptographiquement.
\item
  P0 initial \texttt{P\_total} a été corrigé par addendum v1.0.1 et
  règle P4 \texttt{P\_trainable}; les documents historiques restent
  conservés et explicitement supersédés.
\end{enumerate}

Aucun de ces points ne prouve une intention de sabotage. Un défaut de
traçabilité est un défaut de traçabilité ; il doit être documenté et
corrigé, sans lui attribuer une cause intentionnelle en l'absence
d'éléments.

\subsection{6. Réponse aux critères de
recherche}\label{ruxe9ponse-aux-crituxe8res-de-recherche}

\begin{longtable}[]{@{}
  >{\raggedright\arraybackslash}p{(\columnwidth - 2\tabcolsep) * \real{0.5000}}
  >{\raggedright\arraybackslash}p{(\columnwidth - 2\tabcolsep) * \real{0.5000}}@{}}
\toprule\noalign{}
\begin{minipage}[b]{\linewidth}\raggedright
Critère
\end{minipage} & \begin{minipage}[b]{\linewidth}\raggedright
Résultat EX12
\end{minipage} \\
\midrule\noalign{}
\endhead
\bottomrule\noalign{}
\endlastfoot
Qualité prédictive vs MLP & Mixte ; 11/16 non-infériorités en
interpolation, mais dépendance forte à la tâche. Classification pilote :
log-loss HC-RBF moins bonne. \\
Paramètres à égalité & PASS ; 320 comparaisons P4 et 320 paires P7 à
P\_trainable ±1 \%. GE4X abandonné. \\
Ressources & Échec de l'objectif prioritaire sur CPU mesuré : latence
p50 pipeline batch 1 ≈1,82--2,25× MLP ; stockage C3/C4 plus élevé. \\
Entraînement plus rapide & Non démontré ; ratios de pipeline P5
1,44--1,62× et temps mural P7 \textgreater1× pour les médianes des
quatre variantes. \\
Généralisation & Frontière souvent favorable ; extrapolation défavorable
sur trois des quatre tâches confirmatoires ; shift dépend de la
tâche. \\
\end{longtable}

\subsection{7. Conclusion générale}\label{conclusion-guxe9nuxe9rale}

EX12 ne valide pas l'hypothèse que HC-RBF remplace généralement un MLP
plus efficacement à capacité entraînable égale. L'architecture est
mathématiquement spécifiée, plusieurs variantes ont été réellement
implémentées et comparées, certaines fonctions montrent un avantage
qualitatif, et le mécanisme de couverture apporte des diagnostics
mesurables. Cependant, les coûts de calcul et d'entraînement sont
nettement plus élevés sur la plateforme CPU mesurée, l'extrapolation
reste fragile et la complémentarité des centres n'est prouvée que dans
un petit nombre de cas.

Les acquis principaux sont : un protocole d'évaluation strict, une
architecture HC-RBF complètement spécifiée, un noyau vectorisé
numériquement équivalent, une confrontation à P\_trainable égal, une
campagne multi-graines et un test confirmatoire verrouillé. Les
résultats négatifs et les lacunes de provenance sont conservés et
rapportés. Le verdict est donc \textbf{scientifiquement mixte mais
négatif au regard de l'objectif prioritaire de vitesse/ressources}.

Aucun résultat ne permet d'affirmer une manipulation ou un sabotage des
expériences. Les défauts techniques identifiés (scripts manquants,
hashes historiques absents, compteurs d'audit corrigés, reports de
schéma) sont documentés dans les fichiers d'audit correspondants.

\subsection{8. Carte des livrables et
reproduction}\label{carte-des-livrables-et-reproduction}

\begin{itemize}
\tightlist
\item
  \texttt{prompt\_ex12.json} et
  \texttt{prompt\_ex12\_v1.0.0\_frozen.json} : prompt exécuté et version
  d'origine figée ; addendum de changement :
  \texttt{00\_governance/PROMPT\_CHANGE\_CONTROL\_v1.0.1.json}.
\item
  \texttt{03\_design/HC\_RBF\_FORMAL\_SPEC.md} : formalisation
  architecturale détaillée.
\item
  \texttt{04\_code/} : scripts de génération, entraînement, évaluation,
  benchmark, audit, analyse statistique et rapport.
\item
  \texttt{05\_configs/} : configurations versionnées.
\item
  \texttt{06\_runs/} et \texttt{07\_checkpoints/} : modèles,
  checkpoints, historiques et résultats par run.
\item
  \texttt{08\_benchmarks/} : registres append-only, latences brutes,
  benchmarks et scoring.
\item
  \texttt{09\_analysis/} : métriques recalculées, comparaisons,
  intervalles, tests et contrôles d'intégrité.
\item
  \texttt{10\_reports/} : rapports P0--P7, rapports de lacunes, incident
  \texttt{manifest\_map}, présent rapport et rapport confirmatoire P7.
\item
  \texttt{11\_reproducibility/} : versions et commandes de reproduction.
\item
  \texttt{12\_release/} : manifeste de livraison, checksums et audit de
  complétude.
\end{itemize}

Les commandes exactes et versions de bibliothèques sont décrites dans
\texttt{11\_reproducibility/} et dans les métadonnées des runs. Les
scripts d'entraînement et de scoring ne chargent jamais les sources
historiques. Le paquet de livraison est contrôlé pour ne contenir que
des artefacts EX12 ; \texttt{last\_sources\_3}, les anciens workspaces
EX3--EX11, archives, sauvegardes et holdouts scellés en sont exclus.

\subsection{9. Signature et portée}\label{signature-et-portuxe9e}

\textbf{Rédacteur :} ChatGPT, assistant OpenAI --- rapport généré à
partir des artefacts et registres présents dans \texttt{EX12\_NEXT\_0},
le 4 novembre 2026.

Cette signature est une attribution de rédaction à l'agent. Elle n'est
pas une signature numérique certifiée, une certification de personne
physique, ni une expertise judiciaire. Les SHA-256 et manifestes du
paquet sont fournis séparément pour permettre une vérification technique
indépendante.

\end{document}
